% TOP_PROBLEM: {"id": 30006764, "title": "Tate's rational leading-term Stark conjecture", "category_id": 1, "category": {"id": 1, "name": "number_theory", "display_name": "Number Theory", "description": "Properties of integers, prime numbers, Diophantine equations.", "slug": "number-theory", "order_index": 1, "created_at": "2026-07-31T15:26:25.670Z"}, "status": "open"}
% FIRST_POSTED: 2026-09-27
% !TeX program = lualatex
% Compile twice: lualatex open_problems_500.tex
% All references and prose are embedded; no BibTeX database or external inputs.
\documentclass[11pt,a4paper]{article}
\usepackage[margin=24mm,headheight=14pt]{geometry}
\usepackage{fontspec}
\setmainfont{Latin Modern Roman}
\setsansfont{Latin Modern Sans}
\setmonofont{Latin Modern Mono}
\usepackage{amsmath,amssymb,mathtools}
\usepackage{unicode-math}
\setmathfont{Latin Modern Math}
\usepackage{microtype}
\usepackage{enumitem,xurl,needspace}
\usepackage[dvipsnames]{xcolor}
\usepackage{hyperref}
\hypersetup{unicode=true,colorlinks=true,linkcolor=MidnightBlue,urlcolor=MidnightBlue,
 pdftitle={500 Mathematical Problem Statements},pdfauthor={Source-annotated compilation},bookmarksnumbered=true}
\usepackage{bookmark}
\usepackage{tocloft}
\setlength{\cftsecnumwidth}{3em}
\cftsetpnumwidth{2.5em}
\cftsetrmarg{4em}
\usepackage{titlesec}
\titleformat{\section}{\Large\bfseries\raggedright}{\thesection}{0.7em}{}
\usepackage{fancyhdr}
\pagestyle{fancy}\fancyhf{}
\fancyhead[L]{\small\sffamily Mathematical problem statements}
\fancyhead[R]{\small Source edition 22}
\fancyfoot[C]{\thepage}
\setlength{\parindent}{0pt}
\setlength{\parskip}{5pt plus 1pt}
\setlength{\emergencystretch}{3em}
\setcounter{tocdepth}{1}
\setcounter{secnumdepth}{1}
\setlist[itemize]{leftmargin=3em,itemsep=5pt,topsep=4pt,parsep=0pt}
\allowdisplaybreaks
\newcommand{\N}{\mathbb{N}}
\newcommand{\Z}{\mathbb{Z}}
\newcommand{\Q}{\mathbb{Q}}
\newcommand{\R}{\mathbb{R}}
\newcommand{\C}{\mathbb{C}}
\newcommand{\F}{\mathbb{F}}
\newcommand{\A}{\mathbb{A}}
\newcommand{\PP}{\mathbb P}
\newcommand{\E}{\mathbb{E}}
\newcommand{\eps}{\varepsilon}
\newcommand{\dd}{\,\mathrm d}
\newcommand{\sref}[2]{\hyperlink{source:#1:#2}{[S#2]}}
\newcommand{\eref}[2]{\hyperlink{extra:#1:#2}{[E#2]}}
% Turn the next switch off for a shorter reading copy. The quotations preserve
% every exactTarget field, including qualifications omitted from a short summary.
\newif\ifshowcatalogue
\showcataloguetrue
\newcommand{\cataloguescope}[1]{\ifshowcatalogue
 \paragraph{Exact scope in the supplied catalogue.}
 \begin{quote}\small #1\end{quote}\fi}
% Compile twice with: lualatex 30006764.tex
\numberwithin{equation}{section}
\hypersetup{pdftitle={Research attempt -- problem 30006764},pdfauthor={Research notebook; original compilation retained}}
\fancyhead[L]{\small\sffamily Research notebook}
\fancyhead[R]{\small\sffamily Problem 30006764 | 27 September 2026}
\begin{document}
\setcounter{section}{0}
% ================================================================
% problemId: 30006764
% Canonical problem ID: 30006764
% exactTarget SHA-256: a1c74bb81d780637bf7fb79c82cc52746409a1ed13da7aff0285763d49845e5c
\clearpage
\section*{\texorpdfstring{Tate's rational leading-term Stark conjecture}{Tate's rational leading-term Stark conjecture}}\label{30006764}
\subsection{Definitions and mathematical statement}
For a Galois extension $L/K$ with group $G$, let $S_L$ be the places above $S$ and set
$X=\ker(\Q[S_L]\xrightarrow{\sum}\Q)$ and
$U=\Q\otimes_{\Z}\mathcal O_{L,S}^{\times}$. For a $\Q[G]$-isomorphism $f:X\to U$, the logarithmic map
$\lambda(u)=-\sum_{w\in S_L}\log|u|_w\,w$
uses product-formula normalized absolute values. For an irreducible character $\chi$, take the determinant of $\lambda f$ on
$\operatorname{Hom}_{\C[G]}(V_{\chi^\vee},\C\otimes X)$ to obtain $R_S(\chi,f)$. Let $L_S^*(0,\chi)$ be the first nonzero Laurent coefficient at zero of the Artin $L$-function with the $S$ factors omitted. The target is
\[
 A_S(\chi,f):=\frac{R_S(\chi,f)}{L_S^*(0,\chi)}\in\Q(\chi),
 \qquad\sigma A_S(\chi,f)=A_S(\chi^\sigma,f).
\]
The empty determinant is one; the character dual and the sign in $\lambda$ are part of the chosen normalization.
\par\smallskip\textit{Formulation and concept references:} \sref{30006764}{1}.
\cataloguescope{For every finite Galois extension L/K of number fields with group G, let S be any finite set of places of K containing the archimedean and ramified places. Let X be the augmentation kernel of the free Q-module on places of L above S, U=Q tensor\_Z O\_(L,S)\textasciicircum{}*, and f:X->U any Q[G]-isomorphism. Define lambda:R tensor U->R tensor X by lambda(u)=-sum\_w log|u|\_w w, using product-formula normalized absolute values. For each irreducible complex character chi of G, set R\_S(chi,f)=det\_C(lambda f acting by postcomposition on Hom\_C[G](V\_(dual chi),C tensor X)) and A\_S(chi,f)=R\_S(chi,f)/L\_S\textasciicircum{}*(0,chi), where L\_S\textasciicircum{}* is the first nonzero Laurent coefficient of the S-truncated Artin L-function and an empty determinant is1. Is A\_S(chi,f) in Q(chi), with sigma(A\_S(chi,f))=A\_S(chi\textasciicircum{}sigma,f) for every sigma in Aut(C)?}
\subsection{Short English statement}
After dividing the appropriate logarithmic regulator by an Artin L-function's leading term, is the result algebraic in the character field and compatible with every field automorphism?
% BEGIN RESEARCH ADDITION ATTEMPT-166
\subsection{Research attempt}

\paragraph{Current frontier.}
The rationality and equivariance statement here is Conjecture 1.1 of
\sref{30006764}{1}. Its Remark 1.2 records independence of the choices of $S$
and $f$; Theorem 1.3 treats totally odd characters using Siegel's theorem,
and Remark 1.4 records other established cases, including rational-valued
characters. These statements must be distinguished from the paper's
stronger, integral and prime-specific Stark results. The more recent
Brumer--Stark theorem over $\Z$ concerns the CM/abelian setting, not a
proof of the arbitrary-character target \eref{30006764}{1}, Theorem 1.1.
The calculations below reconstruct normalization reductions explicitly
and test whether class-number information could supply the missing
character-by-character algebraicity. They are not offered as new
versions of the known choice-independence theorem.

\paragraph{Attack route.}
One route is descent of the determinant from a complex representation
to its character field; this succeeds for rational changes of $f$, but
not automatically for the transcendental regulator itself. A second
route removes the extra places one at a time; the required local
factors can be computed exactly. A third tries to recover individual
Stark ratios from products furnished by class-number formulas. The
third route has a genuine loss of information, which is exposed below.

\paragraph{Attempt: changing $f$ without assuming Schur index one.}
For $a\in\operatorname{Aut}_{\Q[G]}X$, let
\begin{equation}
 d_\chi(a)=\det\bigl(a\mid
 \operatorname{Hom}_{\C[G]}(V_{\chi^\vee},\C\otimes X)\bigr).
\end{equation}
For every positive integer $r$, character orthogonality gives
\begin{equation}\label{30006764:trace}
 \operatorname{Tr}(a^r\mid M_\chi)
 =\frac1{|G|}\sum_{g\in G}\chi(g)
   \operatorname{Tr}_{\Q}(a^r g\mid X),\qquad
 M_\chi=\operatorname{Hom}_{\C[G]}(V_{\chi^\vee},\C\otimes X).
\end{equation}
Indeed, decompose $\C\otimes X$ into irreducible representations tensored
with multiplicity spaces. The trace of $a^rg$ is the sum of the
irreducible characters times the traces on those spaces. Pairing with
$\chi(g)$ selects the $\chi^\vee$ summand. All traces on the right other
than $\chi(g)$ are rational. Newton's identities applied to the first
$\dim M_\chi$ power traces therefore show
\begin{equation}
 d_\chi(a)\in\Q(\chi)^\times,\qquad
 \sigma d_\chi(a)=d_{\chi^\sigma}(a).
\end{equation}
This argument does not require $V_\chi$ to have a model over $\Q(\chi)$;
that would be an unjustified Schur-index shortcut.

If $f_2=f_1a$, determinant multiplicativity gives
\begin{equation}
 A_S(\chi,f_2)=d_\chi(a)A_S(\chi,f_1).
\end{equation}
Thus proving the target for one $\Q[G]$-isomorphism proves it for every
one. The empty multiplicity space has trace zero and determinant one,
so that case is covered. This supplies a direct proof of the $f$ part
of the choice-independence statement recorded in \sref{30006764}{1}.

\paragraph{Attempt: an explicit unramified-place transition.}
Let $v\notin S$ be finite. It is unramified, since $S$ already contains
all ramified places. Write $D=\langle F_v\rangle$ for its cyclic
decomposition group, $f_v=|D|$, and $Y_v=\Q[G/D]$. Forgetting the new
place coordinates and taking their valuations give exact sequences
\begin{equation}
 0\to X_S\to X_{S\cup\{v\}}\to Y_v\to0,
 \qquad
 0\to U_S\to U_{S\cup\{v\}}\xrightarrow{\operatorname{ord}}Y_v\to0.
\end{equation}
The first surjection can balance its augmentation using the old places.
The second becomes surjective over $\Q$ because the $S$-ideal class
group is finite. Both split over $\Q[G]$ by Maschke's theorem. Choose
$f'$ compatible with $f$ on the old summands and with the identity on
the two $Y_v$ quotients.

The chosen normalized absolute values give
$-\log|u|_w=\operatorname{ord}_w(u)\log Nw$ and
$\log Nw=f_v\log Nv$. The regulator is consequently block triangular,
and, with $r_\chi=\dim V_\chi^D$, its quotient determinant is
$(f_v\log Nv)^{r_\chi}$. On the analytic side the removed Euler factor is
\begin{equation}
 E_v(s,\chi)=\det(1-F_v(Nv)^{-s}\mid V_\chi).
\end{equation}
Because $F_v$ has finite order, it is semisimple, and
\begin{equation}
 E_v^*(0,\chi)=(\log Nv)^{r_\chi}
  \det(1-F_v\mid V_\chi/V_\chi^D).
\end{equation}
It follows, with precisely the sign and dual conventions in the target,
that
\begin{equation}\label{30006764:place}
 \frac{A_{S\cup\{v\}}(\chi,f')}{A_S(\chi,f)}
 =\frac{f_v^{r_\chi}}
 {\det(1-F_v\mid V_\chi/V_\chi^D)}.
\end{equation}
The denominator is nonzero. Its characteristic coefficients are in
$\Q(\chi)$ by the same power-trace argument, and they transform with
$\chi$ under $\operatorname{Aut}(\C)$. The equality of invariant
multiplicities for a character and its dual follows by averaging over
$D$. Equation \eqref{30006764:place}, together with change of $f$, reduces
all allowed $S$ to the archimedean and ramified places. It reconstructs
the relevant known reduction rather than assuming an Euler factor is
rational just because it has finite-order eigenvalues
\sref{30006764}{1}, Remark 1.2.

\paragraph{Attempt: an exact quadratic normalization test.}
Take $L=\Q(\sqrt5)$, $K=\Q$, $S=\{\infty,5\}$, and
$\varphi=(1+\sqrt5)/2$. Let $w_+,w_-$ be the real places and $w_5$ the
single place over 5. Set
\begin{equation}
 a=w_+-w_-,\qquad b=w_++w_--2w_5,
 \qquad f(a)=\varphi,\quad f(b)=5.
\end{equation}
Units are written multiplicatively on the right, as elements of the
rationalized unit module. Conjugation sends $\varphi$ to
$-\varphi^{-1}$, so $f$ is $G$-equivariant after torsion is killed.
The two elements have independent logarithms and give the required
isomorphism. The logarithm convention yields
\begin{equation}
 \lambda(\varphi)=-(\log\varphi)a,\qquad
 \lambda(5)=-(\log5)b.
\end{equation}
In particular, $|5|_{w_5}=5^{-2}$; replacing this by $5^{-1}$ would
break the product formula and the computation.

For completeness, the Minkowski ideal-class bound is
$\sqrt5/2<2$, so the class number is one. The positive unit
$\varphi$ is fundamental: a unit $1<u<\varphi$ of norm $+1$ would
have integral trace $u+u^{-1}$ strictly between $2$ and $\sqrt5$;
one of norm $-1$ would have trace $u-u^{-1}$ strictly between $0$
and $1$. Both are impossible. Multiplying any positive unit by a
power of $\varphi$ reduces it to this interval or its endpoints.

The analytic class-number formula, in the normalization used in the
Stark source, gives $\zeta_L^*(0)=-(\log\varphi)/2$ and
$\zeta_{\Q}(0)=-1/2$. Hence
\begin{equation}
 L_S^*(0,\chi_5)=\log\varphi,
 \qquad L_S^*(0,1)=-\tfrac12\log5,
 \qquad A_S(\chi_5,f)=-1,\quad A_S(1,f)=2.
\end{equation}
The ramified Euler factor for $\chi_5$ is one. This is a worked known
case, compatible with \sref{30006764}{1}, Remark 1.4, not evidence that the
higher-character descent problem has disappeared.

\paragraph{Stress test: products cannot replace characterwise descent.}
Even all permutation-character products may lose essential information.
For $G=C_3$, the transitive permutation characters are $1$ and
$1+\chi+\bar\chi$. As a model of the information carried by such
products, set
\begin{equation}
 A(1)=1,\qquad A(\chi)=z,\qquad A(\bar\chi)=\bar z,
 \qquad |z|=1,
\end{equation}
where $z$ is transcendental. Such a $z$ exists because the unit circle
is uncountable and the algebraic numbers are countable. Every product
attached to a permutation character is one, but the individual
nontrivial ratios are not algebraic. This is \emph{not} a construction
of Stark ratios for a number field and not a counterexample. It proves
only that the indicated product data, even with complex-conjugation
compatibility, do not logically force individual algebraicity.

Brauer induction with nontrivial one-dimensional characters is stronger
than permutation-character information; invoking it would require the
appropriate Stark information for those characters, not just analytic
class-number formulas for fields. Similarly, the rational power traces
in \eqref{30006764:trace} apply to $a\in\operatorname{End}_{\Q[G]}X$, not to
the logarithmic map itself. Treating the regulator as such a rational
endomorphism would assume the arithmetic descent that remains to be
proved.

\paragraph{Outcome.}
\textbf{Outcome: Conditional reduction.}

The argument explicitly reduces the target to one choice of $f$ and
minimal $S$, without a Schur-index assumption, and verifies all signs
and local norms in a quadratic test. These reductions are already
known in the literature; their proofs here are normalization checks,
not novelty claims. The class-number-product attack has a demonstrated
information deficit. General characterwise algebraicity and full
$\operatorname{Aut}(\C)$-equivariance remain the unresolved input; no
claim about totally odd or abelian cases has been enlarged to the full
catalogue assertion.
% END RESEARCH ADDITION ATTEMPT-166
\subsection{Sources}
\begin{itemize}\raggedright
\item[\textnormal{[S1]}]\hypertarget{source:30006764:1}{} Andreas Nickel, The strong Stark conjecture for totally odd characters; arXiv2106.05619v1 June10,2021; live metadata lists ProcAMS152(2024)147–162.
\url{https://arxiv.org/pdf/2106.05619}.
{\small\textit{Catalogue formulation source.}}
% BEGIN RESEARCH ADDITION SOURCES-166
\item[\textnormal{[E1]}]\hypertarget{extra:30006764:1}{} Samit Dasgupta, Mahesh Kakde, Jesse Silliman and Jiuya Wang, ``The Brumer--Stark Conjecture over $\Z$,'' arXiv:2310.16399 (2023). Theorem 1.1 and its CM/abelian hypotheses.
\url{https://arxiv.org/abs/2310.16399}.
{\small\textit{An adjacent integral theorem; not a solution of arbitrary-character rational Stark.}}
% END RESEARCH ADDITION SOURCES-166
\end{itemize}

\end{document}
